\phantomsection
\color{uniblue}\section*{\centering{\large{ПЕРЕЧЕНЬ СОКРАЩЕНИЙ}\color{white!0}<.>}}
\addcontentsline{toc}{section}{Перечень сокращений}
\color{black}
\begin{longtable}{>{\raggedright\arraybackslash}m{2cm}>{\raggedright\arraybackslash}m{0.5cm}>{\raggedright\arraybackslash}m{20cm}}
\endfirsthead\endhead\endfoot\endlastfoot
ANSI & -- & Американский национальный институт стандартов; \\
GOOSE & -- & Generic Object-Oriented Substation Event; \\
IRF & -- & Internal Relay Fault; \\
PPS & -- & Pulse Per Second; \\
АВ & -- & автоматический выключатель; \\
АВР & -- & автоматическое включение резерва; \\
АСУ~ТП & -- & автоматизированная система управления технологическими процессами; \\
АУ & -- & автоматическое ускорение; \\
АЦП & -- & аналого-цифровой преобразователь; \\
БИ & -- & блок испытательный; \\
БК & -- & блокировка при качаниях; \\
БЛЗШ & -- & блокировка логической защиты шин; \\
БНН & -- & блокировка при неисправности цепей напряжения; \\
БНТ & -- & блокировка при бросках намагничивающего тока; \\
БЧС & -- & блокировка чувствительных ступеней; \\
ВН & -- & высшее напряжение; \\
ВНР & -- & восстановление нормального режима; \\
ВЭ & -- & выдвижной элемент (выкатная тележка, выкатной элемент); \\
ДЗ & -- & дистанционная защита; \\
ЗАВР & -- & запрет автоматического включения резерва; \\
ЗДЗ & -- & защита от дуговых замыканий; \\
ЗМН & -- & защита минимального напряжения; \\
ЗН & -- & заземляющий нож (разъединитель); \\
ЗОП & -- & защита от обрыва провода; \\
ИО & -- & измерительный орган; \\
ИЧМ & -- & интерфейс человек-машина; \\
КА & -- & коммутационный аппарат; \\
КЗ & -- & короткое замыкание; \\
КИ & -- & контроль изоляции; \\
КНН & -- & контроль наличия напряжения; \\
КОН & -- & контроль отсутствия напряжения; \\
КП & -- & контроллер присоединения; \\
КПОН & -- & комбинированный пусковой орган напряжения; \\
КПруж & -- & контроль пружины привода выключателя; \\
КРВ & -- & контроль ресурса выключателя; \\
КС & -- & контроль синхронизма; \\
КСВ & -- & контроль силового выключателя; \\
КЦВ & -- & контроль цепей выключателя; \\
ЛВ & -- & логика включения; \\
ЛЗШ & -- & логическая защита шин; \\
ЛО & -- & логика отключения; \\
МРВ & -- & механический ресурс выключателя; \\
МТЗ & -- & максимальная токовая защита; \\
НКУ & -- & низковольтное комплектное устройство; \\
НН & -- & низшее напряжение; \\
НП & -- & нулевая последовательность; \\
ОБР & -- & оперативная блокировка разъединителей; \\
ОВ & -- & оперативный вывод; \\
ОЗУ & -- & оперативное запоминающее устройство; \\
ОН & -- & орган направления / отключение нагрузки; \\
ОНМ & -- & орган направления мощности; \\
ОП & -- & обратная последовательность; \\
ОПО & -- & общий пусковой орган; \\
ОС & -- & ожидание синхронизма; \\
ОТ & -- & оперативный ток; \\
ОУ & -- & оперативное ускорение; \\
ПДС & -- & преобразователь дискретных сигналов; \\
ПЗУ & -- & постоянное запоминающее устройство; \\
ПО & -- & пусковой орган / программное обеспечение; \\
ПП & -- & прямая последовательность; \\
ПС & -- & предупредительная сигнализация; \\
ПУ & -- & пульт управления; \\
РАС & -- & регистратор аварийных событий; \\
РЗ & -- & релейная защита; \\
РЗА & -- & релейная защита и автоматика; \\
РФК & -- & реле фиксации команд; \\
РЭ & -- & руководство по эксплуатации; \\
СВ & -- & секционный выключатель; \\
СЗЗ & -- & сигнализация замыкания на землю; \\
СС & -- & сборка сигналов; \\
ТЗНП & -- & токовая защита нулевой последовательности; \\
ТК & -- & токовый контроль; \\
ТН & -- & трансформатор напряжения; \\
ТТ & -- & трансформатор тока; \\
ТУ & -- & телеуправление; \\
УВ & -- & управление выключателем / управляющее воздействие; \\
УВЭ & -- & управление выкатным элементом; \\
УЗН & -- & управление заземляющим ножом (разъединителем); \\
УРОВ & -- & устройство резервирования при отказе выключателя; \\
УС & -- & улавливание синхронизма; \\
ФВН & -- & Функция включения на <<холодную>> нагрузку; \\
ФК & -- & функциональная клавиша; \\
ФСУ & -- & функциональная схема устройства; \\
ЦН & -- & цепи напряжения; \\
ЦП & -- & центральный процессор; \\
ЦСВ & -- & цепи сигнализации выключателя; \\
ЧС & -- & чувствительная ступень; \\
ЭМ & -- & электромагнит; \\
ЭМВ & -- & электромагнит включения; \\
ЭМО & -- & электромагнит отключения; \\
ЭМУ & -- & электромагнит управления. \\
\end{longtable}